% =====================================================================
%  Chapter 14 — Step 8
% =====================================================================

\chapter{Step Eight: Automorphy, the Trace Formula, and the Tower}
\label{ch:automorphy}

The argument has to establish that the objects it manipulates are automorphic,
and that the analytic estimates it wants are available. Those are two different
jobs, and the remark's item 16 tries to do both in one sentence.

\begin{quote}
\small
``Relative trace formula, seen in a different half-plane argument via sixfold
cover of $\mathrm{GL}_5$ automorphy of $\mu$ and reducing via
Cauchy--Schwartz to a bound on the norm of a suitably projected Poincar\'e
series, approximate orthogonality of the columns of the integral structure.''
\end{quote}

\section{The relative trace formula and what it is for}The relative trace formula is the transfer principle for inner forms. If
$G$ is an inner form of $G'$, there is a $\mathbb{Z}_p$-Galois spectral
sequence
\begin{equation}
H^*(\mathbb{G}, \mathbb{F}_p) \;\Longrightarrow\;
H^*_{f-is}(G, \mathbb{F}_p) \oplus H^*_{f-c}(G, \mathbb{F}_p),
\end{equation}
whose differentials involve the relative cohomology $H^*(G', \mathbb{F}_p)^{
\mathcal{H}}$, the continuous cohomology of the Weil restriction. In the
arithmetic of Selmer groups this is how one computes: the Selmer group of a
representation over one group is related to the Selmer group over another, via
a map that is an isomorphism on a subspace.

The point for this argument is narrower and specific. The method wants to
produce a large collection of automorphic objects --- the lifts, the periods,
the integral structure --- and it wants the arithmetic of their cohomology to
be under control. The relative trace formula is the device that transfers
control from the group where the analysis is easy (the totally definite
unitary group, where the Selmer groups are known and the periods are
compact) to the group where it is needed.

\begin{quote}
``Seen in a different half-plane argument'': the transfer is not being used
in the usual direction. The remark's phrasing means that the trace formula is
being used to move a statement across the complex half-plane, in a way the
author does not describe further.
\end{quote}

\section{Why the automorphy of mu is the bottleneck}The remark singles out one object: ``automorphy of $\mu$``. Recall from
Chapter~\ref{ch:step4} that $\mu$ is the family of Hecke characters indexing
the stages of the chain. The claim is that $\mu$ is automorphic --- that is,
that the characters used to index the lifts are themselves the characters of
automorphic forms.

The remark gives the mechanism:

\begin{quote}
\small
``$\mu$ is automorphic from $\theta$ on sixfold cover of $\mathrm{GL}_5$ as a
period integral of thetas, quadric has natural interpretation on corresponding
CM e.g.\ $\mathbb{C}/\mathbb{Z}[\omega]$ as Weil, put Poincar\'e bundle on the
edges to get $\theta$s.''
\end{quote}

The structure is recognisable. One wants the characters $\mu$ to be
automorphic. One route to automorphy is to exhibit them as Fourier
coefficients of automorphic forms. The remark's route: a sixfold cover of
$\mathrm{GL}_5$, on which there is a natural structure --- the quadric, which
in the unitary setting corresponds to a CM curve such as
$\mathbb{C}/\mathbb{Z}[\omega]$, viewed through the Weil representation ---
and the Poincar\'e bundle, put on the edges, produces the theta series whose
coefficients are the $\mu$.

So the claim is a construction: build a variety, put a line bundle on it,
take the theta series of that line bundle, and read the characters off its
Fourier coefficients.

\begin{seamnote}{Where this lives in the development}
\code{Development/AutomorphicArithmetic/AutomorphicVarieties/TheoremsMuLinked/}
is the subtree that carries this: \code{GaloisRepresentationTypes/} (including
\code{ReducibleThreeDim.lean} and \code{ConductorAtMapInjectiveCoeff.lean}),
\code{FreeFlag.lean}, \code{Number.lean}. And the deepest input of all is
\code{GlobalTheory.Mem.LiteralAutomorphy.LitN}, the ``literal automorphy''
that \code{Af.Hypotheses\_2} takes as its first field --- see
Chapter~\ref{ch:inputs}.
\end{seamnote}

The sixfold cover of $\mathrm{GL}_5$ deserves a remark of its own. A
sixfold cover of a group is a finite-index subgroup, and a finite-index
subgroup of $\mathrm{GL}_5$ has a very different set of representations from
$\mathrm{GL}_5$ itself. That it is *six*-fold is not incidental: the number
controls the ratio of the relevant indices, and therefore the constant in the
estimates. This is the kind of detail that only shows up in the formalisation,
and it is exactly the kind of detail that an outside reader cannot check
without reading the code.

\section{The Poincar\'e series, and the Cauchy--Schwarz step}The remark continues: ``reducing via Cauchy--Schwartz to a bound on the norm of
a suitably projected Poincar\'e series, approximate orthogonality of the columns
of the integral structure''.

The Poincar\'e series of a unitary group is the automorphic form of the
highest weight type, obtained by summing the matrix coefficients of the
Borel subgroup. Its norm is the quantity that controls the size of the space
of cusp forms, and hence the number of lifts available.

The argument's structure is the standard one:

\begin{enumerate}
\item There is a space of lifts, of dimension $d$ at each stage, and one wants
it to be large.
\item The size of the space is controlled by the norm of the Poincar\'e series,
which is computed by an integral over a quotient, and the integral is
estimated by Cauchy--Schwarz: $\left|\int f g\right|^2 \le \int|f|^2 \int|g|^2$.
\item The two integrals are the norms of the Poincar\'e series and of the
projection, and the second is controlled by the \emph{approximate orthogonality
of the columns of the integral structure}: the columns --- the different
Hecke characters, or the different components of the decomposition --- are
almost orthogonal, so the projection norm is small.
\end{enumerate}

Step 3 is where the averaging of Chapter~\ref{ch:averaging} earns its keep.
The integral structure --- the way the decomposition is glued across the
characters --- is precisely what makes the columns approximately orthogonal,
and that is what makes Cauchy--Schwarz give a good bound rather than a useless
one.

\begin{quote}
\textbf{Where this is formalised.} The column-wise bounds are in
\code{Formalism/SharpRestriction/Divisibility/ChainOperators/PerCol.lean}
(\code{perColSmol}, the scalar compatibility of the per-column map) and
\code{Formalism/SharpRestriction/GrowthBounds/SegmentDegree/PerCol.lean};
the integrality that the columns must respect is the \code{AIntO} of
Chapter~\ref{ch:step2}; and the Poincar\'e-series side of the estimate is in
\code{Development/WeilBound/CoreStatements/AutomorphicKernels/}, the largest
subtree in the Weil-bound development at 932 files.
\end{quote}

\section{The arithmetic growth: the tower again}The last ingredient the remark's item 18 supplies:

\begin{quote}
\small
``Covolumes, theta series' index, congruences between lifts and forms.''
\end{quote}

Three quantities that must agree:

\begin{description}[style=nextline,leftmargin=0pt,font=\sffamily\bfseries\small,
  itemsep=6pt]
\item[\textbf{Covolumes}] The volume of the quotient
$G\backslash \mathrm{U}(E,n)$, which enters every period integral and every
norm. Its computation requires the class number of the order, and hence the
arithmetic of the CM field --- the same arithmetic the Iwasawa tower
supplies.
\item[\textbf{The index of the theta series}] The index of the lattice
generated by the theta series of the Hermitian forms. This is the
quantitative counterpart of the divisibility of Chapter~\ref{ch:step2}: it
says how many of the forms are linearly independent, and hence how many
lifts are genuinely available rather than coincident.
\item[\textbf{Congruences between lifts and forms}] The congruences that
relate the index of one stage's theta series to the next. These are what make
the tower a tower: the congruence at stage $j$ is what makes the lift at
stage $j+1$ well defined.
\end{description}

So the three items are the three things that must be consistent across the
chain: the volumes (the analytic scale), the indices (the algebraic
independent count), and the congruences (the compatibility). Together with
the Iwasawa tower of Chapter~\ref{ch:averaging}, they are what make a chain of
length $m$ exist rather than merely be conceivable.

\section{What step 8 has and has not shown}\begin{enumerate}
\item \textbf{Automorphy of $\mu$ via a sixfold cover of $\mathrm{GL}_5$.} A
construction: a variety, a line bundle, the theta series, and the characters
read off as Fourier coefficients. In the literature, automorphy of characters
in this generality is exactly the sort of thing that is known in special cases
and open in general; the development's \code{TheoremsMuLinked} subtree is
where the work is done.
\item \textbf{The Poincar\'e-series bound via Cauchy--Schwarz.} Standard
analytic technology, applied to the projected space.
\item \textbf{Approximate orthogonality of the columns.} The property the
integral structure must have, and the reason the averaging of
Chapter~\ref{ch:averaging} is needed at all.
\item \textbf{Volumes, indices, congruences.} The three compatibility
conditions across the chain.
\end{enumerate}

The honest note is this: items 1 and 3 are where the argument is most exposed,
and they are the two places where the remark's own vocabulary is least
informative. ``Sixfold cover of $\mathrm{GL}_5$'' and ``congruences between
lifts and forms'' are not descriptions that let a reader check anything. They
are pointers.

\begin{keytakeaways}
\item The relative trace formula transfers control of the cohomology from the
group where the analysis is easy to the group where it is needed, in a
direction the remark calls ``a different half-plane argument''.
\item Automorphy of the characters $\mu$ is obtained by exhibiting them as
Fourier coefficients of theta series on a sixfold cover of $\mathrm{GL}_5$,
with the quadric interpreted on a CM curve.
\item The analytic estimate is a Cauchy--Schwarz bound on the norm of a
projected Poincar\'e series, and it works only because the columns of the
integral structure are approximately orthogonal.
\item Volumes, theta indices and congruences are the three quantities that must
agree across the chain; they are where the Iwasawa tower enters.
\item The two least transparent items --- the sixfold cover and the
congruences --- are also the two most load-bearing, and they are where an
outside reader is most dependent on the code.
\end{keytakeaways}
